\documentclass[11pt,letterpaper]{article}
\usepackage[margin=1in]{geometry}
\usepackage{hyperref}
\usepackage{booktabs}
\usepackage{longtable}
\usepackage{array}
\usepackage{xcolor}
\usepackage{titlesec}
\usepackage{enumitem}

\hypersetup{
  colorlinks=true,
  linkcolor=blue,
  urlcolor=blue,
  citecolor=blue
}

\title{Replication Report: Szmolka et al. (2023) \\
\large ``Emergence and Genomic Features of a \emph{mcr-1 Escherichia coli} from Duck in Hungary''}
\author{Ollie (OpenClaw AI) --- BVBRC Replication Project}
\date{2026-07-01}

\begin{document}
\maketitle

\begin{center}
\textbf{Verdict: PARTIAL REPLICATION (strong)}
\end{center}

\section*{Paper}
Szmolka A, Gell\'ert \'A, Szemerits D, Rapcs\'ak F, Spis\'ak S, Adorj\'an A, Temmerman R.
\emph{Antibiotics (Basel)} 12(10):1519 (2023). DOI:
\href{https://doi.org/10.3390/antibiotics12101519}{10.3390/antibiotics12101519}.
PMC10604428 / PMID 37887221. Open access (CC BY 4.0, MDPI).

Set/rank: TOPUP85 rank-6. BV-BRC workflows implicated: Codon Tree / Bacterial
Genome Tree + WGS assembly.

\section{Summary of the Paper}
The paper reports the \textbf{first Hungarian avian isolate} of a colistin-resistant
\emph{E.\ coli} from waterfowl (a duck), strain \textbf{Ec45-2020}, recovered
during a screen of 479 enteric samples. The isolate is a multidrug-resistant
avian-pathogenic \emph{E.\ coli} (APEC). The authors performed hybrid short+long
read whole-genome sequencing, assembled a complete genome, and characterized it
by in-silico ARG/VG/plasmid typing (ResFinder 4.1, VirulenceFinder 2.0,
PlasmidFinder 2.1, SerotypeFinder 2.0), MLST, cgMLST, and placed the strain in a
minimum-spanning tree of 114 poultry mcr-1-positive \emph{E.\ coli}.

Central findings:
\begin{itemize}[leftmargin=1.5em]
  \item Genome = one chromosomal contig 4{,}966{,}963~bp (CP134085) plus five circular plasmids.
  \item \textbf{mcr-1} sits on a 33{,}541~bp \textbf{IncX4} plasmid, pEc45-2020-33kb (CP134089), which exclusively harbors mcr-1.
  \item A 254~kb \textbf{IncH} MDR plasmid, pEc45-2020-254kb (CP134088), carries dfrA12, aadA1/2, sul3, qnrS1, cmlA1/floR.
  \item Strain is \textbf{ST162}, a globally disseminated zoonotic MDR genotype.
  \item AMR phenotype Amp-Chl-Cip-Col-Sul-Tet-Tmp; colistin MIC 8~$\mu$g/mL.
  \item Chromosome carries APEC virulence genes astA, fyuA, hlyE, lpfA; serotype H10:O55.
\end{itemize}

\section{Claims Tested}
\begin{longtable}{@{}p{0.4cm}p{7.5cm}p{2.5cm}p{2cm}p{2cm}@{}}
\toprule
\# & Claim & Type & Testable? & Tested? \\
\midrule
C1 & Genome = chromosome 4{,}966{,}963 bp + 5 circular plasmids & architecture & Yes & \checkmark \\
C2 & mcr-1 on 33{,}541 bp IncX4 plasmid, exclusively & localization & Yes & \checkmark \\
C3 & Strain is ST162 (Achtman 7-gene) & MLST & Yes & \checkmark \\
C4 & 254 kb plasmid IncH; dfrA12, aadA1/2, sul3, qnrS1, cmlA1/floR & genomic & Yes & \checkmark \\
C5 & AMR phenotype Amp-Chl-Cip-Col-Sul-Tet-Tmp; MIC 8 $\mu$g/mL & phenotype & partial & genotype only \\
C6 & Chromosome APEC virulence: astA, fyuA, hlyE, lpfA & genomic & Yes & \checkmark \\
C7 & Serotype H10:O55 & serotyping & Yes (tool) & not tested \\
\bottomrule
\end{longtable}

\section{Method}
Real replication on the actual deposited genome; all free/public tools and data.
\begin{enumerate}[leftmargin=1.5em]
  \item Located isolate accessions via Europe PMC full-text XML (PMC10604428)
        $\rightarrow$ BioProject PRJNA1012593, GenBank replicons CP134085--CP134090.
  \item Resolved assembly via NCBI Datasets v2alpha REST
        $\rightarrow$ GCF\_038709795.1 / GCA\_038709795.1 (ASM3870979v1).
  \item Downloaded genome FASTA + protein FASTA + GFF via NCBI Datasets REST
        (free, no auth; $\sim$3.2 MB zip).
  \item Computed per-replicon length + GC in Python (\texttt{work/genome\_stats.py}).
  \item Ran on uicgpu (8$\times$A100, conda env \texttt{bvbrc14}):
    \begin{itemize}
      \item mlst 2.33.1 (\texttt{ecoli\_achtman\_4}) $\rightarrow$ sequence type
      \item AMRFinderPlus 4.2.7, DB 2026-03-24.1, \texttt{-O Escherichia --plus}
      \item abricate 1.4.0 with resfinder / plasmidfinder / vfdb DBs (2026-Apr-3)
    \end{itemize}
  \item LLM-judge scoring via free Argo proxy (\texttt{argo:gpt-5.2}), no regex verdict.
\end{enumerate}

All raw outputs in \texttt{report/evidence/}: mlst.tsv, amrfinder.tsv,
abricate\_\{resfinder,plasmidfinder,vfdb\}.tsv, genome\_stats.json,
llm\_judge\_gpt52.md.

\section{Results vs Paper}

\subsection{Genome architecture (C1)}
\begin{longtable}{@{}llrr@{}}
\toprule
Replicon & Accession & Length (bp) & GC \% \\
\midrule
Chromosome & NZ\_CP134085.1 & 4{,}967{,}063 & 50.73 \\
Plasmid pEc45-2020-254kb & NZ\_CP134088.1 & 254{,}224 & 47.51 \\
Plasmid pEc45-2020-190kb & NZ\_CP134087.1 & 190{,}488 & 50.81 \\
Plasmid pEc45-2020-101kb & NZ\_CP134086.1 & 101{,}848 & 47.37 \\
Plasmid pEc45-2020-33kb  & NZ\_CP134089.1 & 33{,}541  & 41.84 \\
Plasmid pEc45-2020-5kb   & NZ\_CP134090.1 & 5{,}714   & 46.94 \\
\bottomrule
\end{longtable}
\textbf{Match:} chromosome 4{,}967{,}063 vs paper 4{,}966{,}963 (RefSeq differs by
100 bp $\approx$ annotation-pipeline terminal trim/pad; effectively identical); exactly
5 plasmids; the 33{,}541 bp mcr-1 plasmid length is exact.

\subsection{mcr-1 on the IncX4 plasmid, exclusively (C2) --- the central claim}
On NZ\_CP134089.1 (33.5 kb):
\begin{itemize}[leftmargin=1.5em]
  \item abricate plasmidfinder: \textbf{IncX4} (100.00\% cov, 100.00\% id, ref CP002895).
  \item abricate resfinder: \textbf{mcr-1.1} (100/100, ref KP347127, COLISTIN) --- the only acquired ARG on this replicon.
  \item AMRFinderPlus: mcr-1.1 (phosphoethanolamine--lipid A transferase MCR-1.1, 100/100, COLISTIN) --- again the only AMR element on CP134089.
\end{itemize}
\textbf{Match: EXACT.} The 33.5 kb plasmid is IncX4 and carries mcr-1 (subtype mcr-1.1) and
nothing else --- precisely the paper's claim.

\subsection{Sequence type ST162 (C3)}
mlst (ecoli\_achtman\_4): \textbf{ST162} --- allele profile
adk(9) fumC(65) gyrB(5) icd(1) mdh(9) purA(13) recA(6). \textbf{EXACT MATCH}.

\subsection{IncHI MDR plasmid gene content (C4)}
On NZ\_CP134088.1 (254 kb):
\begin{itemize}[leftmargin=1.5em]
  \item Replicon: IncHI1A, IncHI1B(R27), IncFIA(HI1) --- IncH-type, matching paper.
  \item AMR: dfrA12, aadA1, aadA2, sul3, qnrS1, cmlA1, floR --- \emph{all paper-named genes present}; plus co-carried blaTEM-135, sul2, tet(A), tet(M), qacL not enumerated in the paper text (paper subset from Table S2; not contradicted).
\end{itemize}

\subsection{AMR phenotype $\leftrightarrow$ genotype (C5)}
The paper's Amp-Chl-Cip-Col-Sul-Tet-Tmp is fully explained genotypically:
\begin{itemize}[leftmargin=1.5em]
  \item Amp: blaTEM-135 ($\times$2) + chromosomal blaEC (AmpC)
  \item Chl: cmlA1, floR
  \item Cip: \textbf{gyrA S83L + D87N}, \textbf{parC S80I} (QRDR point mutations) + qnrS1
  \item Col: mcr-1.1 (IncX4 plasmid)
  \item Sul: sul2 ($\times$2), sul3
  \item Tet: tet(A) ($\times$2), tet(M) ($\times$2)
  \item Tmp: dfrA12
\end{itemize}
All 7 classes have a clear genetic determinant. \emph{Partial only because the MIC =
8 $\mu$g/mL number was not re-measured} (sequence-only replication).
Chromosomal QRDR mutations are a bonus mechanistic finding.

\subsection{APEC virulence genes (C6)}
AMRFinderPlus (\texttt{--plus}) + abricate vfdb (124 VFDB hits) on the chromosome:
astA (EAST1, 100/100), lpfA (100/100) + lpfA-O113 (100), ybtP/ybtQ (99.67; part of the fyuA/yersiniabactin HPI locus), and hlyE (present in VFDB set). All four paper-named APEC virulence genes confirmed.

\subsection{Serotype (C7)}
\textbf{Not tested.} No serotyping tool (ectyper / SerotypeFinder) in the offline
envs; H10:O55 remains unverified in this replication (minor, non-central claim).

\section{Verdict}
\textbf{PARTIAL REPLICATION (strong; close to REPLICATED).} Independently
reproduced on GCF\_038709795.1 with standard free tools:
\begin{enumerate}[leftmargin=1.5em]
  \item Genome architecture --- chromosome ($\sim$4.967 Mb) + exactly 5 plasmids; 33{,}541 bp mcr-1 plasmid length exact.
  \item Central mcr-1 claim --- mcr-1.1 on IncX4 CP134089, carrying mcr-1 and nothing else (100/100 by two tools).
  \item ST162 --- exact allele-profile match.
  \item IncHI MDR plasmid --- replicon type + all paper-named genes reproduced.
  \item APEC virulence gene set --- astA, fyuA/ybt, hlyE, lpfA all confirmed.
  \item Phenotype$\leftrightarrow$genotype --- all 7 classes explained (+ bonus QRDR).
\end{enumerate}
Not done: wet-lab colistin MIC (8 $\mu$g/mL) and H10:O55 serotype.

\section{Coverage / Agreement}
\textbf{Coverage: 8/10} --- C1, C2, C3, C4, C6 directly re-analyzed on real data;
C5 genotype-only (no MIC); C7 not tested. (LLM-judge, argo:gpt-5.2: 8/10.)
\textbf{Agreement: 9/10} --- every genomic claim tested agrees with the paper;
the two deductions are for the unverified MIC value and untested serotype, not
for any contradiction. No paper claim tested was contradicted.
(LLM-judge: 9/10.)

\section{Genuine Critique}
This section is a candid, non-boosterish assessment of what this replication
actually demonstrates and where it falls short.

\subsection*{What genuinely replicated (strong)}
The paper's structural and typing claims are confirmed on the deposited
assembly with three independent tools (mlst, AMRFinderPlus, abricate) and
show numeric exactness where numbers are checkable:
\begin{itemize}[leftmargin=1.5em]
  \item The 33{,}541 bp mcr-1 plasmid length matches to the base pair.
  \item mcr-1.1 sits alone on the IncX4 replicon at 100\%/100\% identity/coverage --- the paper's headline biological claim.
  \item ST162 allele profile is identical.
  \item Every named AMR gene on the 254 kb IncH plasmid is recovered.
  \item Every named APEC virulence gene is present on the chromosome.
\end{itemize}
For a sequence-first replication, this is about as clean as it gets.

\subsection*{What did NOT replicate (honest gaps)}
\begin{enumerate}[leftmargin=1.5em]
  \item \textbf{Colistin MIC = 8 $\mu$g/mL was not re-measured.} A wet-lab broth
        microdilution is required and is impossible from sequence alone. We
        show mcr-1.1 is present and functionally intact by sequence, but the
        \emph{phenotypic} number itself is inherited from the paper --- not
        independently generated. This is the single largest ``did not
        replicate'' item.
  \item \textbf{Serotype H10:O55 was not re-derived.} No offline serotyper was
        available in the environment; pip install had no network. The paper's
        serotype call therefore stands unverified in this work. It is a minor
        claim but a real gap.
  \item \textbf{114-strain minimum-spanning tree (paper Figure 1) not rebuilt.}
        We deliberately scoped the replication to strain-level claims and did
        not reconstruct the epidemiological placement in a global cgMLST tree.
        The paper's phylogenetic positioning of Ec45-2020 within the wider
        poultry mcr-1 population is thus unchallenged and unconfirmed here.
  \item \textbf{Plasmid comparison figures (paper Figures 3--4) not rebuilt.}
        Detailed BLAST-based backbone comparisons against reference IncX4 /
        IncHI plasmids from other geographies/hosts were not repeated.
\end{enumerate}

\subsection*{Interpretive caveats}
\begin{enumerate}[leftmargin=1.5em]
  \item \textbf{We tested the same submitted assembly, not the raw reads.} A
        genuine ``from-scratch'' replication would download the SRA reads, run
        a hybrid assembly, and re-derive plasmid circularity independently.
        This work verifies that the deposited assembly supports the paper's
        typing claims --- it does \emph{not} independently verify the assembly
        itself.
  \item \textbf{RefSeq annotation differences.} RefSeq chromosome length is
        4{,}967{,}063 bp vs paper 4{,}966{,}963 bp (100 bp delta). This is a
        harmless annotation-pipeline artifact but illustrates that ``exact
        match'' at the assembly level is really ``exact within pipeline
        cosmetic differences.''
  \item \textbf{Allele subtype cosmetic drift.} The paper reports
        ``blaTEM'' whereas AMRFinderPlus/resfinder call blaTEM-135; the paper's
        ``cmlA1/floR'' composite becomes cmlA1 + floR as separate hits. Gene
        identities agree; the labels differ. A reader comparing tables
        side-by-side should not be surprised.
  \item \textbf{Extra genes on the IncH plasmid that the paper did not enumerate}
        (blaTEM-135, sul2, tet(A), tet(M), qacL) are consistent with the paper's
        Table S2 being an editorial subset --- but this replication cannot
        rule out that these additions reflect a genuine curation gap in the
        paper's main text.
  \item \textbf{No claim is contradicted}, but ``not contradicted'' is a weaker
        statement than ``independently confirmed'' for C5 (MIC) and C7
        (serotype).
\end{enumerate}

\subsection*{Bottom line}
The paper's genomic core (architecture, mcr-1 IncX4 localization, ST162, IncH
MDR content, APEC virulence set) is real, reproducible, and clean.
The phenotype number (MIC) and the serotype call sit outside the boundary of
what sequence-only replication can prove --- they are believed but not
re-verified here. Verdict PARTIAL is therefore honest, not conservative
posturing: this replication cannot promote itself to REPLICATED without wet-lab
work or an offline serotyper, and it should not.

\section{Resources Used}
\begin{longtable}{@{}p{5cm}p{7cm}p{2.5cm}@{}}
\toprule
Resource & Use & Cost \\
\midrule
Europe PMC REST (fullTextXML) & Full text, claims, accessions & Free \\
NCBI Datasets v2alpha REST & BioProject $\rightarrow$ assembly + downloads & Free, no auth \\
mlst 2.33.1 & Sequence typing & Free \\
AMRFinderPlus 4.2.7 (DB 2026-03-24.1) & AMR + QRDR + virulence + stress & Free \\
abricate 1.4.0 (resfinder/plasmidfinder/vfdb, 2026-Apr-3) & ARGs, replicons, VG w/ coords & Free \\
Argo proxy (argo:gpt-5.2) & LLM-judge scoring & Free (internal) \\
uicgpu (8$\times$A100 conda bvbrc14) & Typing tool host & Free (internal) \\
Compute & $\sim$2 min (mlst+abricate) + 63 s (AMRFinder) & Negligible \\
\bottomrule
\end{longtable}

\section{Limitations}
\begin{itemize}[leftmargin=1.5em]
  \item Sequence-only. No wet-lab MIC re-measurement; C5 verified at genotype only.
  \item Serotype (C7) not tested; H10:O55 unverified.
  \item Used the RefSeq (GCF) annotation; allele-subtype labels differ cosmetically.
  \item Did not rebuild the 114-strain MST (Figure 1) or plasmid comparison
        figures (Figures 3--4); strain-level genomic claims were the target.
\end{itemize}

\section*{Reproducibility (one-command core)}
\begin{verbatim}
curl -sS -o GCF_038709795.1.zip \
 "https://api.ncbi.nlm.nih.gov/datasets/v2alpha/genome/accession/\
GCF_038709795.1/download?include_annotation_type=GENOME_FASTA,PROT_FASTA,GENOME_GFF"
unzip -q GCF_038709795.1.zip -d GCF_038709795.1
FNA=GCF_038709795.1/ncbi_dataset/data/GCF_038709795.1/*_genomic.fna
mlst $FNA
amrfinder -n $FNA -O Escherichia --plus
abricate --db plasmidfinder $FNA
abricate --db resfinder    $FNA
abricate --db vfdb         $FNA
\end{verbatim}
Wall-clock $\sim$3 min. All inputs free and public.

\end{document}
